\subsubsection{Liquid Metal Breeders modeling}

To assess technologies for tritium breeding and cooling of high heat flux fusion components, we propose to conduct
never-before-performed high-fidelity simulations of (i) breeder blanket modules and (ii) divertor. In particular, we propose to expand the work we performed on the WCLL blanket and perform coupled simulations involving MHD, tritium extraction, and heta transfer at realistic conditions. We also seek to perform divertor simulations.  

The tasks for the WCLL breeder module are:

\begin{enumerate}
\item \textbf{Complete LES and MHD of WCLL module}. We will also quantify appropriateness of RANS for liquid metal MHD in blankets based on bulk transport properties (e.g. effective conductivity).
\item\textbf{Coupling to Monte Carlo and Tritium transport codes}. The goals is to understand how local/global coupling influences peak heat flux and other metrics. Extract information on bounding operation conditions.
\end{enumerate}

The tasks for the divertor simulations are:

\begin{enumerate}
\item \textbf{Build baseline model} of a realistic divertor component, including CAD meshing for radiation transport and hexahedral meshes for spectral element CFD. Multiphysics coupling will be included after first establishing mesh convergence and robustness of each single physics.
\item \textbf{High-fidelity multiphysics analysis} of the divertor, and quantify appropriateness of various combinations of single- or partial-physics and whether coarse mesh models can be formulated to replicate full-detail solutions. 
\end{enumerate}


